\section{Raíces}\label{sec:raices}

Los conectores dicen \emph{cómo} se relacionan las ideas; las raíces dicen
\emph{de qué} hablo. Con veintinueve raíces nuevas y ocho signos que ya existían
se pueden formar cientos de conceptos, del mismo modo que el chino o los
Blissymbolics combinan unas pocas piezas.

La meta no es dejar de usar el español. Los nombres propios y los términos
técnicos (Hume, salitre, NP) se siguen escribiendo con letras, igual que en
matemática se escribe «Gauss». La meta es que las ideas generales que se
repiten en todas las materias (mente, poder, vida, cambio, verdad\ldots) tengan
signo, y que un apunte quede a mitad de camino: más o menos la mitad en
KalaScript y la mitad en español.

\subsection{Las raíces}
\textbf{Mente y persona}\par\vspace{1mm}
\KalaTarjetas{R01}{R02}{R03}{R04}
\KalaTarjetas{R05}{R12}{R21}{R22}
\textbf{Vida y naturaleza}\par\vspace{1mm}
\KalaTarjetas{R06}{R07}{R08}{R28}
\KalaTarjetas{R15}{R16}{R17}{R26}
\textbf{Sociedad y valor}\par\vspace{1mm}
\KalaTarjetas{R09}{R10}{R11}{R25}
\noindent\KalaTarjeta{R27}\hfill\KalaTarjetaVacia\hfill\KalaTarjetaVacia\hfill\KalaTarjetaVacia\par\vspace{3mm}
\textbf{Forma, técnica y percepción}\par\vspace{1mm}
\KalaTarjetas{R13}{R14}{R20}{R23}
\KalaTarjetas{R24}{R19}{R18}{R29}

\subsection{Ocho signos que también son raíces}
Estos signos ya existían. Además de su uso habitual, funcionan como raíces
dentro de un compuesto:
\begin{center}\renewcommand{\arraystretch}{1.5}
\begin{tabular}{@{}cl@{\qquad}cl@{}}
\kss{K30} & persona & \kss{K08} & lugar \\
\kss{K07} & tiempo & \kss{K01} & cosa \\
\kss{K05} & cantidad & \kss{K04} & cambio \\
\kss{K06} & grupo & \kss{K19} & lo desconocido, la pregunta \\
\end{tabular}
\end{center}

\subsection{Cómo se forma un compuesto}
\begin{enumerate}
\item \textbf{Primero el núcleo.} La primera raíz dice qué \emph{es} la cosa; las
 siguientes precisan de qué tipo. Grupo \(+\) persona es un grupo de personas:
 \cps{K06}{K30} «sociedad». Persona \(+\) poder es una persona con poder:
 \cps{K30}{R09} «gobernante».
\item \textbf{Un lazo debajo.} Las raíces de un compuesto se escriben juntas y
 se unen con una sola línea por debajo. Sin el lazo son palabras separadas.
\item \textbf{Dos raíces, como mucho tres.} Si hacen falta más, es mejor
 escribir la palabra o definir una abreviatura.
\item \textbf{Es una etiqueta, no una definición.} \cps{K01}{R01} «cosa de la
 mente» se lee «idea», aunque no todo lo mental sea una idea. El diccionario
 fija la lectura.
\item \textbf{Si no es obvio, se declara.} Un compuesto que no está en el
 diccionario se escribe arriba de la página con su palabra la primera vez, o
 se añade a la leyenda personal.
\end{enumerate}

\subsection{Raíces sueltas y con otros signos}
\begin{itemize}
\item Una raíz sola es una palabra: \kss{R07} «vida», \kss{R09} «poder».
\item Raíces y compuestos sirven de sustantivo o de verbo según el lugar:
 \kss{R19} es «ver» o «lo visto»; \cps{R02}{K19} es «duda» o «dudar».
\item Con \(\neg\) se niega: \(\neg\)\kss{R07} es «no vivo», que \emph{no} es lo
 mismo que \kss{R08} «muerto»: una piedra no está viva, pero tampoco muerta.
\item Con \(\uparrow\) y \(\downarrow\): \(\uparrow\)\cps{K05}{R10} «sube el precio».
\item \textbf{Conflicto entre dos nombres} se lee «contra»: «Chile \kss{R25} Perú».
\item \textbf{Origen entre dos cosas} se lee «viene de»: «idea \kss{R26} impresión».
\end{itemize}

\clearpage
\subsection{Diccionario}\label{sec:diccionario}
\newcommand{\Voz}[2]{#1 & #2 \\}
\noindent
\begin{minipage}[t]{.32\linewidth}\small\renewcommand{\arraystretch}{1.5}
\textbf{Materias} \(=\) saber \(+\) su objeto\par\vspace{1mm}
\begin{tabular}{@{}l@{\quad}l@{}}
\Voz{\cps{R02}{R13}}{matemática}
\Voz{\cps{R02}{R15}}{física}
\Voz{\cps{R02}{R16}}{química}
\Voz{\cps{R02}{R07}}{biología}
\Voz{\cps{R02}{R06}}{medicina}
\Voz{\cps{R02}{R01}}{psicología}
\Voz{\cps{R02}{K07}}{historia}
\Voz{\cps{R02}{R10}}{economía}
\Voz{\cps{R02}{R11}}{derecho}
\Voz{\cps{R02}{R03}}{lingüística}
\Voz{\cps{R02}{R24}}{computación}
\Voz{\cps{R02}{R12}}{filosofía}
\Voz{\cps{R02}{R21}}{ética}
\Voz{\cps{R02}{R28}}{geografía}
\Voz{\cps{R02}{R19}}{ciencia (empírica)}
\end{tabular}\par\vspace{3mm}
\end{minipage}\hfill
\begin{minipage}[t]{.32\linewidth}\small\renewcommand{\arraystretch}{1.5}
\textbf{Mente}\par\vspace{1mm}
\begin{tabular}{@{}l@{\quad}l@{}}
\Voz{\cps{K01}{R01}}{idea}
\Voz{\cps{R01}{K07}}{memoria}
\Voz{\cps{R01}{R11}}{razón}
\Voz{\cps{R02}{R04}}{experiencia}
\Voz{\cps{R02}{K19}}{duda}
\Voz{\cps{K04}{R02}}{aprendizaje}
\Voz{\cps{R02}{R14}}{teoría}
\Voz{\cps{R01}{R03}}{significado}
\Voz{\cps{R01}{R05}}{voluntad}
\Voz{\cps{R01}{R19}}{percepción}
\Voz{\cps{R06}{R04}}{sensación}
\Voz{\cps{R04}{R21}}{placer}
\Voz{\cps{R04}{R22}}{dolor, sufrimiento}
\Voz{\cps{R20}{K07}}{hábito}
\Voz{\cps{R01}{R27}}{alma}
\end{tabular}\par\vspace{3mm}
\end{minipage}\hfill
\begin{minipage}[t]{.32\linewidth}\small\renewcommand{\arraystretch}{1.5}
\textbf{Sociedad y economía}\par\vspace{1mm}
\begin{tabular}{@{}l@{\quad}l@{}}
\Voz{\cps{K06}{K30}}{sociedad}
\Voz{\cps{K06}{R09}}{gobierno}
\Voz{\cps{K30}{R09}}{gobernante}
\Voz{\cps{K08}{R09}}{país, Estado}
\Voz{\cps{K08}{K06}}{ciudad}
\Voz{\cps{R25}{K06}}{guerra}
\Voz{\cps{R25}{K08}}{batalla}
\Voz{\cps{K04}{R09}}{revolución}
\Voz{\cps{R03}{R11}}{tratado, acuerdo}
\Voz{\cps{R21}{R11}}{justicia}
\Voz{\cps{R20}{R22}}{delito}
\Voz{\cps{K06}{R27}}{religión}
\Voz{\cps{K08}{R10}}{mercado}
\Voz{\cps{K05}{R10}}{precio}
\Voz{\cps{K06}{R10}}{empresa}
\Voz{\cps{R20}{R10}}{trabajo (empleo)}
\end{tabular}\par\vspace{3mm}
\end{minipage}

\vspace{4mm}\noindent
\begin{minipage}[t]{.32\linewidth}\small\renewcommand{\arraystretch}{1.5}
\textbf{Vida}\par\vspace{1mm}
\begin{tabular}{@{}l@{\quad}l@{}}
\Voz{\cps{R06}{R07}}{organismo}
\Voz{\cps{K01}{R07}}{célula}
\Voz{\cps{K06}{R07}}{especie}
\Voz{\cps{K04}{R07}}{evolución}
\Voz{\cps{R26}{R07}}{nacimiento}
\Voz{\cps{K08}{R07}}{ecosistema}
\Voz{\cps{R06}{R21}}{salud}
\Voz{\cps{R06}{R22}}{enfermedad}
\end{tabular}\par\vspace{3mm}
\textbf{Naturaleza y materia}\par\vspace{1mm}
\begin{tabular}{@{}l@{\quad}l@{}}
\Voz{\cps{K08}{R28}}{la Tierra, planeta}
\Voz{\cps{K07}{R28}}{clima}
\Voz{\cps{K04}{R16}}{reacción química}
\Voz{\cps{K05}{R15}}{temperatura}
\Voz{\cps{R09}{R17}}{fuerza}
\Voz{\cps{K05}{R17}}{velocidad}
\Voz{\cpp{K05}{K04}{R17}}{aceleración}
\Voz{\cps{R14}{R17}}{onda}
\end{tabular}\par\vspace{3mm}
\end{minipage}\hfill
\begin{minipage}[t]{.32\linewidth}\small\renewcommand{\arraystretch}{1.5}
\textbf{Forma y verdad}\par\vspace{1mm}
\begin{tabular}{@{}l@{\quad}l@{}}
\Voz{\cps{R11}{R13}}{función}
\Voz{\cps{K04}{R12}}{demostración}
\Voz{\cps{R26}{R12}}{axioma}
\Voz{\cps{R25}{R12}}{contradicción}
\Voz{\cps{R12}{R13}}{ecuación, igualdad}
\end{tabular}\par\vspace{3mm}
\textbf{Información y técnica}\par\vspace{1mm}
\begin{tabular}{@{}l@{\quad}l@{}}
\Voz{\cps{R23}{R24}}{computadora}
\Voz{\cps{R11}{R20}}{algoritmo}
\Voz{\cps{R11}{R23}}{programa}
\Voz{\cps{K08}{R24}}{memoria (de máquina)}
\Voz{\cps{K06}{R23}}{red}
\Voz{\cps{R03}{R23}}{lenguaje de programación}
\end{tabular}\par\vspace{3mm}
\end{minipage}\hfill
\begin{minipage}[t]{.32\linewidth}\small\renewcommand{\arraystretch}{1.5}
\textbf{Lenguaje y arte}\par\vspace{1mm}
\begin{tabular}{@{}l@{\quad}l@{}}
\Voz{\cps{K01}{R03}}{palabra}
\Voz{\cps{R11}{R03}}{gramática}
\Voz{\cps{R03}{R29}}{literatura}
\Voz{\cps{R20}{R29}}{arte}
\Voz{\cps{R19}{R29}}{pintura}
\Voz{\cps{R18}{R29}}{música}
\end{tabular}\par\vspace{3mm}
\end{minipage}

Fíjate en que el núcleo cambia el significado: \cps{R01}{K07} (mente con tiempo)
es la memoria de una persona, y \cps{K08}{R24} (lugar de información) es la de
una máquina.

\needspace{14\baselineskip}
\subsection{Cómo se trazan}
Casi todas las raíces son de uno o dos trazos. Las que más se parecen entre sí
se distinguen por su interior:
\begin{center}
\begin{tabular}{@{}cp{.62\linewidth}@{}}
\toprule Par & Qué los distingue \\ \midrule
\KalaPar{R01}{R02} & Mente es una cúpula abierta con un punto; Saber, la cúpula cerrada y vacía. \\[3mm]
\KalaPar{R21}{R22} & Bien lleva una cruz dentro; Mal, solo la barra. \\[3mm]
\KalaPar{R26}{K27} & Origen tiene un punto en el centro; Parte de, un círculo pequeño a un lado. \\[3mm]
\KalaPar{R07}{R08} & Vida tiene la hoja hacia arriba; Muerte, la hoja caída. \\[3mm]
\KalaPar{R04}{R05} & Sentir es un corazón; Querer, un corazón pequeño con un hilo. \\[3mm]
\KalaPar{R23}{K34} & Máquina tiene una rueda dentro de la caja; Revisar está vacía. \\[3mm]
\KalaPar{R25}{K25} & Conflicto son dos puntas rectas y una barra; Pero, dos arcos. \\
\bottomrule
\end{tabular}
\end{center}
